\documentclass[UTF8,a4paper,10pt]{ctexart}
\usepackage{latexsym}
\usepackage[empty]{fullpage}
\usepackage{titlesec}
\usepackage{marvosym}
\usepackage[usenames,dvipsnames]{color}
\usepackage{verbatim}
\usepackage{enumitem}
\usepackage[colorlinks=true,urlcolor=blue,linkcolor=blue]{hyperref}
\usepackage{fancyhdr}
\usepackage{tabularx}
\usepackage{microtype}

% Adjust margins to fit summary and new content on one page

\addtolength{\oddsidemargin}{-0.75in}
\addtolength{\evensidemargin}{-0.75in}
\addtolength{\textwidth}{1.5in}
\addtolength{\topmargin}{-0.7in}
\addtolength{\textheight}{1.4in}

\urlstyle{same}
\raggedbottom
\raggedright
\setlength{\tabcolsep}{0in}

% Section formatting
\titleformat{\section}{\vspace{-7pt}\heiti\bfseries\raggedright\large}{}{0em}{}[\color{black}\titlerule \vspace{-7pt}]

% Macros
\newcommand{\resumeItem}[1]{\item\small{\linespread{1.1}\selectfont #1\par}\vspace{-0.5pt}}
\newcommand{\resumeSubheading}[4]{%
  \vspace{-2pt}\item
  \begin{tabular*}{0.97\textwidth}[t]{l@{\extracolsep{\fill}}r}
    \textbf{#1} & #2 \\
    #3 & #4 \\
  \end{tabular*}\vspace{-7pt}
}
\newcommand{\resumeProjectHeading}[2]{%
  \item
  \begin{tabular*}{0.97\textwidth}{l@{\extracolsep{\fill}}r}
    \normalsize #1 & #2 \\
  \end{tabular*}\vspace{-4pt}
}
\newcommand{\resumeSubHeadingListStart}{\begin{itemize}[label={}, leftmargin=0.15in]}
\newcommand{\resumeSubHeadingListEnd}{\end{itemize}}
\newcommand{\resumeItemListStart}{\begin{itemize}[label=\tiny$\bullet$, itemsep=0pt, topsep=3pt, parsep=0pt]}
\newcommand{\resumeItemListEnd}{\end{itemize}\vspace{-6pt}}

\begin{document}

% HEADER
\begin{center}
  \textbf{\large \heiti 陈以珊 (Yishan Chen)} \\ \vspace{1pt}
  \small AI Automation \& Data Analysis $|$ Multi-Agent Systems $|$ Workflow Automation \\ \vspace{1pt}
  \small \href{mailto:Kristyn921207@gmail.com}{Kristyn921207@gmail.com} $|$ 微信: kristyn5942 $|$ \href{https://kristynchen.github.io/Personal-Website/}{\underline{个人网站}} $|$ \href{https://github.com/kristynchen}{\underline{GitHub}}
\end{center}
\vspace{-12pt}

% 个人简介
\section{个人简介}
\small {电子与计算机工程专业学生，具备 Python 自动化、多智能体系统 (Agentic AI) 与数据分析实战经验。主导构建 UAT 自动化测试框架与 n8n 工作流管线，擅长利用 AI 工具优化业务流程并提升运营效率。具备优秀英语听说读写能力与多国交换经历，善于跨文化团队协作与沟通。}
\vspace{-4pt}

% 教育背景
\section{教育背景}
\resumeSubHeadingListStart
  \resumeProjectHeading
    {\textbf{上海交通大学} \quad \quad | \quad 密西根学院 电子与计算机工程学士 \quad | \quad 中国·上海 }{2026年8月（预计毕业）}
    \resumeItemListStart
      \resumeItem{\textbf{奖学金}：台湾学生奖学金（2023--2025）；\textbf{百贤亚洲未来领袖奖学金 AFLSP }——面向亚洲顶尖高校学生的跨文化交流及领导力项目。}
      \resumeItem{\textbf{相关课程}：STAT4060J Computational Methods, STAT4710J Data Science and Analytics using Python}
    \resumeItemListEnd

  \resumeProjectHeading
    {\textbf{庆应义塾大学} \quad \quad | \quad 计算机科学（交换生） \quad | \quad 日本·东京}{2025年3月 -- 2025年8月}
  \resumeProjectHeading
    {\textbf{特鲁瓦科技大学} \quad | \quad 法语（交换生） \quad | \quad 法国·特鲁瓦}{2025年1月 -- 2025年2月}
  \resumeProjectHeading
    {\textbf{关西学院大学} \quad \quad | \quad 法律（交换生） \quad | \quad 日本·兵库}{2024年1月 -- 2024年2月}
\resumeSubHeadingListEnd

% 工作与教学经历
\section{工作与教学经历}
\resumeSubHeadingListStart
  \resumeProjectHeading
    {\textbf{上海比特无限智能科技有限公司} \quad | \quad AI 应用工程师实习生}{2026年3月 -- 至今}
    \resumeItemListStart
      \resumeItem{主导构建自动化 UAT 验收测试框架：针对 130+ 工业指令场景建立多维度评分标准，通过优化 Prompt 逻辑将系统意图识别准确率从 27\% 提升至 86\%，具备丰富的 UAT 测试设计与执行落地经验。}
      \resumeItem{实现数字化看板自动化管线：打通测试数据与 Obsidian 看板的闭环联动，实现 KPI 指标实时可视化与自动报表输出。}
      \resumeItem{探索 \textbf{AI + 工作流优化 (AI Innovation)}：利用 VLM (视觉大模型) 实现结构化数据提取与异常识别，将繁杂的文档审核转化为自动化流水线，成功减少约 \textbf{90\% 的人工覆核成本}，符合团队对 AI 隐患识别的技术预研需求。}
    \resumeItemListEnd

  \resumeProjectHeading
    {\textbf{上海交通大学} \quad | \quad ENRG1000J 工程导论 助教}{2023年9月 -- 2025年12月}
    \resumeItemListStart
      %\resumeItem{开发 Python 自动化脚本，实现作业批量处理与异常数据识别，显著减少人工操作时间，优化课程运营流程；具备独立完成数据处理闭环的实战能力。}
      \resumeItem{主动识别助教工作中的重复性痛点，自行构建 n8n 自动化 Pipeline（Gmail $\times$ Google Sheets），覆盖学生评量解析、合格证书自动发送、改分记录汇总及期末作业分级打包交付全流程；每学期节省 8+ 小时人工操作，消除漏发/错发风险；工作流设计具备高度可复用性。}
      \resumeItem{带领 60+ 名本地与国际混合班学生完成每周 Lab 课程；\textbf{连续三次获评「优秀助教」}并受邀留任，展现跨文化沟通与多元团队协作能力。}
    \resumeItemListEnd
\resumeSubHeadingListEnd

% 项目经历
\section{项目经历}
\resumeSubHeadingListStart

  \resumeProjectHeading
    {\textbf{\href{https://terny.ai/}{terny.ai} — AI 科研机会匹配系统} \quad | \quad Python, LangGraph, LLM}{2025年11月 -- 至今}
    \resumeItemListStart
      \resumeItem{构建基于爬虫与大语言模型的科研信息采集系统，实现复杂实验室网站与科研动态的自动化解析与结构化抽取。}
      \resumeItem{设计并实现基于 RAG 的多级检索策略（Excel 权威数据优先 + Web fallback + 向量检索），将高校识别准确率从 75\% 提升至 98\%+。}
      \resumeItem{主导 Agentic AI 多智能体架构设计与落地，构建协同决策流水线，实现研究方向理解、用户建模与双向匹配，并通过事件日志与执行追踪提升系统可观测性与调试效率。}
    \resumeItemListEnd

  \resumeProjectHeading
    {\textbf{假新闻文本分类与数据可视化 Dashboard} \quad | \quad Python, R, Shiny}{2025年9月 -- 2025年12月}
    \resumeItemListStart
      \resumeItem{清洗处理 \textbf{7 万条}新闻数据，使用 TF-IDF + SVM 模型实现文本分类，\textbf{准确率达 93\%}，优于基准模型。}
      \resumeItem{使用 Shiny 构建\textbf{交互式可视化看板}，支持模型预测实时监测与数据动态查询，直接对应业务运营可视化需求。}
    \resumeItemListEnd

  \resumeProjectHeading
    {\textbf{房价预测模型与数据分析} \quad | \quad Python, Excel, scikit-learn}{2025年7月 -- 2025年9月}
    \resumeItemListStart
      \resumeItem{完成数据清洗、特征工程与多模型集成，实现约 \textbf{0.36 RMSE}，排名班级\textbf{前 2 名}；使用 Excel 制作统计表辅助结果分析，熟练运用数据透视表与公式进行数据整理与报表呈现。}
    \resumeItemListEnd

  \resumeProjectHeading
    {\textbf{基于 GNN 与 Transformer 的材料性质预测引擎}\quad | \quad Python, PyTorch, GNNs, Transformer}{2025年12月 -- 至今}
    \resumeItemListStart
      \resumeItem{主导自动化数据处理流，数据涵盖材料学级（OC22/DFT）样本清洗聚合、特征抽离及多维度节点向量化映射计算。}
      \resumeItem{搭建并在代码维度实现基于图神经网络（GNNs）与 Transformer 的计算底座，高度重构并客制化具备跨模型可复用的核心 PyTorch 计算集群代码。}
    \resumeItemListEnd

  %\resumeProjectHeading
    %{\textbf{基于 GNN 与 Transformer 的材料性质预测引擎}}{2025年12月 -- 至今}
    %\resumeItemListStart
      %\resumeItem{主导自动化数据处理流，数据涵盖材料学级（OC22/DFT）样本清洗聚合、特征抽离及多维度节点向量化映射计算。}
      %\resumeItem{搭建并在代码维度实现基于图神经网络（GNNs）与 Transformer 的计算底座，高度重构并客制化具备跨模型可复用的核心 PyTorch 计算集群代码。}
      %\resumeItem{自研并封装完整的业务指标检测及训练过程管线网络，在确保特征层与张量演变更迭具备内部\textbf{可观测性}基础上，为大规模云测与持续工业部署提供可用技术切口。}
    %\resumeItemListEnd
\resumeSubHeadingListEnd

% 领导力与活动
\section{领导力与活动}
\resumeSubHeadingListStart
  \resumeProjectHeading
    {\textbf{台北市立第一女子高级中学学生会副主席 — E-campus 数字化校园项目负责人}}{2020年12月 -- 2022年6月}
    \resumeItemListStart
      \resumeItem{主导校园数字化转型项目，联合台北市政府推行线上点餐与自助设备；项目落地后学生平均排队时间缩短 80\%，热食部单日营业额提升 100\%，日均服务 800+ 学生，累计受益学生逾 6000 人。}
      \resumeItem{独立识别学生痛点并完成跨部门提案，协调学校、政府、供应商三方完成系统部署，并代表团队出席新闻发布会进行项目汇报。}
    \resumeItemListEnd

  \resumeProjectHeading
    {\textbf{上海交通大学密西根学院青年志愿者 — 新闻宣传部副部长}}{2022年9月 -- 2024年8月}
    \resumeItemListStart
      \resumeItem{主导品牌宣传与社群运营：通过精确的内容推送策略，在 \textbf{3 天内成功吸引 60+ 位}报名者参与志愿活动，显著提升组织影响力。}
      \resumeItem{带领 30 人团队赴异地开展两周支教，负责项目整体调度与跨多元背景的沟通协调，优化团队协作效率。}
    \resumeItemListEnd
\resumeSubHeadingListEnd

% 技能
\section{技能}
\vspace{2pt}
\small
\textbf{编程语言}：Python, R, SQL, C/C++, MATLAB \\
\textbf{数据分析与可视化}：Excel（数据透视表、公式、图表）、pandas, NumPy, scikit-learn, Shiny Dashboard, 数据清洗 \\
\textbf{AI 与自动化}：LangGraph, VLM, LLM Prompt Engineering, n8n, AI 辅助脚本开发, AI 工作流优化 \\
\textbf{办公工具}：Excel, PowerPoint, Outlook, Git, LaTeX \\
\textbf{语言能力}：中文（母语），英语（流利），日语（会话）

\end{document}
